\section{结论}

我们开发了一个用于序列数据概率建模的 u-MPS 模型，并证明它在并行性、采样与 regularization 方面具有独特的能力，其方式呼应了正则语言（regular languages）的结构。虽然我们的结果是在 u-MPS 这一特定语境下推导的，但用于展示这些能力的基本技术同样适用于具有类似数学结构的其他模型，例如 WFA、HMM 与 PSR。因此，我们的定理~\ref{thm:sample} 与~\ref{thm:regularization} 应当可以推广到其他此类模型。我们预期本文开发的算法在应用于其他模型时会伴随不同的运行时间，从而使并行化、采样与 regularization 方法呈现与本文所报告的不同的性能权衡。例如，这里使用的并行求值方法对 u-MPS 需要 $\bigO(n D^3)$ 的资源，而对 HQMM 则需要 $\bigO(n D^6)$。

除了这些直接的推广之外，我们结果的一个更有趣的扩展是将定理~\ref{thm:sample} 与~\ref{thm:regularization} 推广到上下文无关语言（context-free languages）的情形。虽然所有正则语言都是上下文无关的，但后者具有显著更强的表达能力，并且与自然语言处理和自动代码生成中的应用更为相关。令人惊讶的是，我们发现这样的语言确实可以在 u-MPS 模型中被采样并用作 regularizer，尽管代价更高，为 $\bigO(D^6)$。

鉴于我们展示了利用文法（grammar）来约束 u-MPS 模型概率分布的能力，一个自然的问题是逆过程是否可行：即给定一个训练好的 u-MPS 模型，是否存在自动的手段来识别能够解释所学概率分布中部分相关性的文法规则？这样的技术将代表一类定性上全新的文法推断（grammatical inference），其中文法规则与语言模型以双向方式相互作用。

最后，一个自然的下一步是将 u-MPS 扩展到真实世界的序列建模任务，尤其是语言建模。当前这一过程的障碍有：(a) 某些 u-MPS 操作 $\bigO(D^3)$ 的代价，以及 (b) 缺乏用梯度下降训练大型张量网络的成熟最佳实践。我们预期这些障碍将被专门的工程投入以及数量迅速增长的处理张量网络的软件库所克服，同时多体物理社区开发的强大计算方法也将被机器学习所采纳。鉴于这里展示出的意外收益，我们期望循环张量网络架构拥有光明的未来。

\section*{致谢}
本研究得到加拿大高等研究院（CIFAR AI chair program）的支持。

\appendix
\onecolumn

